% Figure: the 555-timer astable oscillator schematic (timing network R_A,
% R_B, C) on the left and its capacitor and output waveforms on the right.
% Plain TikZ for the schematic; pgfplots for the waveforms.
\begin{figure}[htbp]
\centering
% === left: timing network schematic ===
\begin{tikzpicture}[
  line width=0.8pt,
  res/.style={draw, engblue, fill=white, minimum width=4mm, minimum height=9mm, font=\scriptsize},
  font=\small,
]
  % 555 block
  \draw[engblue] (2.4,0.0) rectangle (4.4,3.0);
  \node[font=\scriptsize] at (3.4,2.6) {555};
  \node[font=\tiny, anchor=east] at (2.4,2.4) {DIS};
  \node[font=\tiny, anchor=east] at (2.4,1.5) {THR};
  \node[font=\tiny, anchor=east] at (2.4,0.6) {TRG};
  \node[font=\tiny, anchor=west] at (4.4,1.5) {OUT};
  % Vcc rail
  \draw[engred] (0.0,4.0) -- (5.0,4.0);
  \node[engred, font=\scriptsize, anchor=south] at (2.5,4.0) {$V_{CC}$};
  \draw[engred] (3.4,3.0) -- (3.4,4.0);
  % R_A from Vcc to DIS
  \node[res] (RA) at (0.6,3.2) {$R_A$};
  \draw[engblue] (0.6,4.0) -- (RA);
  \draw[engblue] (0.6,2.75) -- (0.6,2.4) -- (2.4,2.4);
  % R_B from DIS to THR/TRG node
  \node[res] (RB) at (0.6,1.6) {$R_B$};
  \draw[engblue] (0.6,2.75) -- (RB);
  \draw[engblue] (0.6,1.15) -- (0.6,0.6);
  % node joining THR and TRG
  \draw[engblue] (0.6,1.5) -- (2.4,1.5);
  \draw[engblue] (0.6,0.6) -- (2.4,0.6);
  \node[circle, fill=engblack, inner sep=0.7pt] at (0.6,1.5) {};
  % capacitor C from THR/TRG node to ground
  \draw[engblue] (1.5,1.5) -- (1.5,0.95);
  \node[circle, fill=engblack, inner sep=0.7pt] at (1.5,1.5) {};
  \draw[engblue, line width=1.1pt] (1.25,0.95) -- (1.75,0.95);
  \draw[engblue, line width=1.1pt] (1.25,0.80) -- (1.75,0.80);
  \node[font=\scriptsize, anchor=west] at (1.8,0.87) {$C$};
  \draw[engblue] (1.5,0.80) -- (1.5,0.3) -- (3.4,0.3);
  \node[font=\tiny] at (3.4,0.05) {gnd};
  % output
  \draw[engblue] (4.4,1.5) -- (5.0,1.5);
  \node[font=\scriptsize, anchor=west] at (5.0,1.5) {$v_o$};
\end{tikzpicture}
\hspace{0.3cm}
% === right: waveforms ===
\begin{tikzpicture}
\begin{axis}[
  width=6.6cm, height=5.0cm,
  xlabel={$t/T$}, ylabel={voltage},
  xmin=0, xmax=2, ymin=-0.1, ymax=1.25,
  axis lines=left,
  tick label style={font=\tiny},
  label style={font=\scriptsize},
  ytick={0.333,0.667,1.0},
  yticklabels={$\tfrac{1}{3}V_{CC}$,$\tfrac{2}{3}V_{CC}$,$V_{CC}$},
  yticklabel style={font=\tiny},
  legend style={font=\tiny, at={(0.98,0.30)}, anchor=east, draw=none, fill=none},
  clip=false,
]
% capacitor voltage: charges 1/3 -> 2/3 over fraction 0.64 of period, then
% discharges 2/3 -> 1/3 over 0.36 of period. Sketch as piecewise exponentials.
% Charge phase 0 .. 0.64 : 0.333 + (1 - 0.333) approached; use shaped curve to 0.667
\addplot[engblue, very thick, domain=0:0.64, samples=80]
  {1 - (1 - 0.333)*exp(-1.72*x/0.64)};
\addplot[engblue, very thick, domain=0.64:1.0, samples=60]
  {0.667*exp(-1.10*(x-0.64)/0.36)};
\addplot[engblue, very thick, domain=1.0:1.64, samples=80]
  {1 - (1 - 0.333)*exp(-1.72*(x-1.0)/0.64)};
\addplot[engblue, very thick, domain=1.64:2.0, samples=60]
  {0.667*exp(-1.10*(x-1.64)/0.36)};
\addlegendentry{$v_C$}
% output: high during charge (0 .. 0.64), low during discharge
\addplot[engorange, thick] coordinates {
  (0,1.1) (0.64,1.1) (0.64,-0.05) (1.0,-0.05) (1.0,1.1) (1.64,1.1)
  (1.64,-0.05) (2.0,-0.05)};
\addlegendentry{$v_o$}
\draw[enggray, dashed] (axis cs:0,0.333) -- (axis cs:2,0.333);
\draw[enggray, dashed] (axis cs:0,0.667) -- (axis cs:2,0.667);
\end{axis}
\end{tikzpicture}
\caption[555-timer astable oscillator and its waveforms]{The 555-timer
astable oscillator (left) and its waveforms (right). The capacitor
charges toward $V_{CC}$ through $R_A+R_B$ until it reaches the upper
threshold $\tfrac{2}{3}V_{CC}$, at which point the discharge transistor
turns on and the capacitor discharges through $R_B$ alone toward the
lower threshold $\tfrac{1}{3}V_{CC}$, where the cycle restarts. The
output is high during the charge phase and low during the discharge
phase, so the asymmetry between the charge path ($R_A+R_B$) and the
discharge path ($R_B$) sets a duty cycle always above one half.}
\label{fig:vol04ch09:555-astable}
\end{figure}
